\section{Representations of valued quivers\\ and their Hall--Ringel
algebras}\label{sec:valued quivers}

Fix a~f\/inite f\/ield $\mathbb{F}$ and an algebraic closure $\bar{\mathbb{F}}$.
For each positive integer~$d$ write $\mathbb{F}_d$ for the degree~$d$ extension of $\mathbb{F}$ in $\bar{\mathbb{F}}$.
Note that with these conventions we may intersect $\mathbb{F}_d$ and $\mathbb{F}_{d'}$ to get that
$\mathbb{F}_{\gcd(d,d')}$ is their largest common subf\/ield.
We now def\/ine representations of $(Q,\mathbf{d})$ by assigning an $\mathbb{F}_{d_i}$-vector space to each vertex $i\in
I$ and an $\mathbb{F}_{\gcd(d_i,d_j)}$-linear map to each arrow from vertex~$i$ to vertex~$j$.
The f\/inite-dimensional representations of $(Q,\mathbf{d})$ form a~f\/initary, hereditary, Abelian category
$\operatorname{rep}_\mathbb{F}(Q,\mathbf{d})$ where kernels and cokernels are taken vertex-wise, we refer the reader  %, where
to~\cite{rupel1} for more details.

Recall that we work under the assumption that the quiver~$Q$ contains no oriented cycles.
Then the simple representations of $(Q,\mathbf{d})$ may be labeled as $S_i$ ($i\in I$) where the only non-zero vector  %, where
space is $\mathbb{F}_{d_i}$ at vertex~$i$.
The Grothendieck group $\mathcal{K}(Q,\mathbf{d})$ of $\operatorname{rep}_\mathbb{F}(Q,\mathbf{d})$ has
a~$\mathbb{Z}$-basis given by the classes $\alpha_i=[S_i]$ ($i\in I$) of the irreducible representations $S_i$, in this
way we identify the root lattice $\mathcal{Q}$ with $\mathcal{K}(Q,\mathbf{d})$.
Write $[V]$ for the isomorphism class of a~representation $V\in\operatorname{rep}_\mathbb{F}(Q,\mathbf{d})$ and write
$|V|$ for its \emph{dimension vector}, i.e.~the class of~$V$ in $\mathcal{K}(Q,\mathbf{d})$.
Def\/ine the \emph{height} $\operatorname{ht}(V)$ of a~representation~$V$ as the length of any composition series for~$V$.

We def\/ine the \emph{Euler--Ringel form}
$\langle\cdot,\cdot\rangle:
\mathcal{K}(Q,\mathbf{d})\times\mathcal{K}(Q,\mathbf{d})\to\mathbb{Z}$ on generators~by
\begin{gather*}
\langle\alpha_i,\alpha_j\rangle=
\begin{cases}
d_i & \text{if $i=j$,}
\\
-[d_ib_{ij}]_+ & \text{if $i\ne j$.}
\end{cases}
\end{gather*}
The identif\/ication of $\mathcal{Q}$ and $\mathcal{K}(Q,\mathbf{d})$ allows to transfer the bilinear form $(\cdot,\cdot)$
to $\mathcal{K}(Q,\mathbf{d})$.
It then follows from the def\/initions that $(\cdot,\cdot)$ is the symmetrization of the Euler--Ringel form, %--
i.e.~$(\alpha_i,\alpha_j)=\langle\alpha_i,\alpha_j\rangle+\langle\alpha_j,\alpha_i\rangle$ for all $i,j\in I$.
Recall that we write $\alpha_i^\vee=\alpha_i/d_i$.
For a~dimension vector $\mathbf{v}\in\mathcal{K}(Q,\mathbf{d})$ def\/ine ${}^*\mathbf{v}\in\mathcal{K}(Q,\mathbf{d})$~by
\begin{gather*}
{}^*\mathbf{v}=\sum\limits_{i\in I}\langle\alpha_i^\vee,\mathbf{v}\rangle\alpha_i.
\end{gather*}
